\documentclass[11pt]{article}
\usepackage[margin=1in]{geometry}
\usepackage{amsmath,amssymb,amsthm}
\usepackage{xurl}
\usepackage{hyperref}
\sloppy
\let\oldus\_ \renewcommand{\_}{\oldus\allowbreak}
\newtheorem{theorem}{Theorem}
\newtheorem{lemma}[theorem]{Lemma}
\theoremstyle{definition}
\newtheorem{definition}[theorem]{Definition}
\newcommand{\Reff}{R_{\mathrm{eff}}}

\title{Proof of conjecture 00000003419:\\ commute time $=2m\cdot$ effective resistance}
\author{Jackmeson1}
\date{2026-10-05}

\begin{document}
\maketitle

\section{The conjecture and its reading}

\textbf{Statement (English, verbatim).} Definition: Hitting times: the symmetry of commute times.
Conjecture: The commute-time symmetry: the two-point commute time is twice the edge count times the
resistance, with the proof of symmetry the symmetric difference of Green functions. (commute
resistance twice symmetric difference)

The Chinese text says the same; its Definition line contains the stray word ``Holant'', which plays
no role in the claim and is ignored.

\textbf{Reading.} $G=(V,E)$ is a finite, connected, simple (unweighted) graph with $m=|E|$ edges.
The walk is the simple random walk; every edge is a resistor of one ohm. For vertices $a,b$ the
claim is
\[ \kappa(a,b) := H(a,b)+H(b,a) = 2m\,\Reff(a,b), \]
together with the symmetry $\kappa(a,b)=\kappa(b,a)$ and $\Reff(a,b)=\Reff(b,a)$. The clause ``the
proof of symmetry [is] the symmetric difference of Green functions'' describes a proof method. We
follow it through the hitting-time vectors $h_b=H(\cdot,b)$: the unit-current potential is the
normalized difference $(h_b-h_a)/(2m)$ (Lemma~\ref{lem:pot}). Green functions are not built in Lean
as separate objects. The bridge to the Green-function formulation is as follows: with
$\Gamma=L^{+}$ (the Green function of the Laplacian on the orthogonal complement of the constants),
$L\,\Gamma(\cdot,a)=e_a-\tfrac1n\mathbf 1$, so for $a\ne b$ the vector $\Gamma(\cdot,a)-\Gamma(\cdot,b)$
solves $Lv=e_a-e_b$; since the kernel of $L$ is the constants (connectedness), our established
equation $L\big((h_b-h_a)/(2m)\big)=e_a-e_b$ gives
$(h_b-h_a)/(2m)=\Gamma(\cdot,a)-\Gamma(\cdot,b)$ modulo an additive constant. Thus the symmetric
difference of Green functions is exactly the potential used in the proof (this remark is not
formalized; the main identity does not depend on it).

\textbf{Conventions.} $T_b=\min\{t\ge 0: X_t=b\}$, so $H(b,b)=0$. Connectedness is needed, since
otherwise hitting times are infinite. The case $a=b$ is allowed, and both sides are then $0$.
Weighted graphs and multigraphs are not treated.

\textbf{Attribution.} The identity is due to A.~K. Chandra, P. Raghavan, W.~L. Ruzzo and
R. Smolensky, \emph{The electrical resistance of a graph captures its commute and cover times},
first presented at STOC 1989; the journal version, with P. Tiwari as an additional author, appeared in
Computational Complexity 6 (1996) 312--340 (Crossref record retrieved:
\url{https://doi.org/10.1007/bf01270385}). Wikipedia's article
``Resistance distance'' (\url{https://en.wikipedia.org/wiki/Resistance_distance}) defines the resistance distance of a simple connected graph as the resistance
``with each edge being replaced by a resistance of one ohm''. It describes the commute time as ``the
expected number of steps in a random walk that starts at'' $u$, visits $v$ and returns to $u$, and
states $C_{u,v}=2m\Omega_{u,v}$ with that reference. Our proof is self-contained and fully
formalized.

\section{Definitions (as in the Lean file)}

Let $L=D-A$ be the Laplacian (Mathlib \texttt{SimpleGraph.lapMatrix}).
\begin{definition}
\begin{enumerate}
\item Transition probability: $P(x,y)=1/\deg x$ if $x\sim y$, else $0$ (\texttt{trans}). Lemma
\texttt{sum\_trans}: $\sum_y P(x,y)=1$ when $\deg x>0$.
\item Trajectory probability of $w=(w_0,\dots,w_t)$: $\prod_{i<t}P(w_i,w_{i+1})$
(\texttt{pathProb}). On graphs where every vertex has positive degree this is the law of the
Markov chain on the first $t+1$ steps. On the connected one-vertex graph, \texttt{trans} is
identically $0$, so \texttt{pathProb} is not a probability law for $t>0$; there the only hitting
target is the starting vertex, every survival probability is $0$, and both sides of the identity
vanish, so the theorem still holds (equivalently, one may adopt an absorbing self-loop convention
at an isolated vertex).
\item $P_x(T_b>t)=\sum\{\prod_{i<t}P(w_i,w_{i+1}) : w_0=x,\ w_i\ne b \text{ for all } i\le t\}$
(\texttt{survProb}).
\item $H(x,b)=E_x[T_b]=\sum_{t\ge 0}P_x(T_b>t)$, the tail-sum formula for the expectation of a
random variable with values in $\{0,1,2,\dots\}\cup\{\infty\}$ (\texttt{hittingTime}, a real
\texttt{tsum}; the series is proved to converge). $\kappa(a,b)=H(a,b)+H(b,a)$
(\texttt{commuteTime}).
\item A unit-current potential from $a$ to $b$ is a $v:V\to\mathbb R$ with $Lv=e_a-e_b$
(\texttt{IsUnitPotential}): one unit of current enters at $a$ and leaves at $b$, and Kirchhoff's
and Ohm's laws hold with unit resistors. $\Reff(a,b)=v(a)-v(b)$ (\texttt{effRes}, chosen by
\texttt{Classical.choose}). We prove that such a $v$ exists and that $v(a)-v(b)$ does not depend on
the choice. Equivalently, $\Reff(a,b)=(e_a-e_b)^{\top}L^{+}(e_a-e_b)$ with $L^+$ the Moore--Penrose
inverse; this remark is not formalized.
\end{enumerate}
\end{definition}

\section{Proof}

\begin{lemma}[\texttt{survProb\_zero}, \texttt{survProb\_succ}, \texttt{survProb\_self}]
$P_x(T_b>0)=[x\ne b]$, $P_b(T_b>t)=0$, and
$P_x(T_b>t+1)=[x\ne b]\sum_y P(x,y)P_y(T_b>t)$.
\end{lemma}
\begin{proof}
Split a trajectory of length $t+1$ as $(w_0,w')$ using \texttt{Fin.consEquiv}. The product
factors as $P(w_0,w'_0)\cdot\prod_{i<t} P(w'_i,w'_{i+1})$. Then sum over $w_0=x$ and group by $w'_0=y$.
\end{proof}

\begin{lemma}[\texttt{exists\_dirichlet}, \texttt{dirichlet\_nonneg}]
There is $g$ with $g(b)=0$, $(Lg)(x)=\deg x$ for $x\neq b$, and $g\ge 0$.
\end{lemma}
\begin{proof}
Let $M$ be $L$ with row $b$ replaced by $e_b^{\top}$. If $Mg=0$, then $g(b)=0$ and $(Lg)(x)=0$ for
$x\ne b$. Column sums of $L$ vanish (\texttt{sum\_lapMatrix\_mulVec}, from $L^{\top}=L$ and $L\mathbf
1=0$), so $(Lg)(b)=0$ too. On a connected graph $\ker L$ consists of the constants
(\texttt{lapMatrix\_mulVec\_eq\_zero\_iff\_forall\_reachable}), so $g=0$. Thus $M$ is injective,
hence invertible and surjective, and we solve $Mg=(\deg x\,[x\ne b])_x$. Minimum principle: let
$g$ attain its minimum at $x_0$. If $x_0\ne b$, then $\deg x_0>0$ by connectedness, and
$\deg x_0=(Lg)(x_0)=\sum_{y\sim x_0}(g(x_0)-g(y))\le 0$, a contradiction. So the minimum is
$g(b)=0$.
\end{proof}

\begin{lemma}[\texttt{partial\_le}, \texttt{summable\_survProb}]
$\sum_{t<N}P_x(T_b>t)\le g(x)$ for all $N$, so $\sum_t P_x(T_b>t)$ converges.
\end{lemma}
\begin{proof}
For $x\ne b$, $(Lg)(x)=\deg x$ means $g(x)=1+\sum_yP(x,y)g(y)$ (\texttt{lap\_eq\_degree\_iff}).
Induct on $N$. Using the recursion and $P\ge0$, we get
$\sum_{t<N+1}P_x(T_b>t)=1+\sum_yP(x,y)\sum_{t<N}P_y(T_b>t)\le 1+\sum_yP(x,y)g(y)=g(x)$. The
terms are nonnegative, so \texttt{summable\_of\_sum\_range\_le} applies.
\end{proof}

\begin{lemma}[\texttt{hittingTime\_step}, \texttt{lap\_hittingTime}]
$H(b,b)=0$, and $H(x,b)=1+\sum_yP(x,y)H(y,b)$ for $x\ne b$. Hence
$(Lh_b)(x)=\deg x-2m[x=b]$.
\end{lemma}
\begin{proof}
Shift the convergent series by one and exchange the finite sum with the series. For $x\neq b$ the
first-step equation is $(Lh_b)(x)=\deg x$. At $b$, vanishing column sums and
$\sum_x\deg x=2m$ (\texttt{sum\_degrees\_eq\_twice\_card\_edges}) give
$(Lh_b)(b)=-\sum_{x\ne b}\deg x=\deg b-2m$.
\end{proof}

\begin{lemma}[\texttt{hitting\_potential}]\label{lem:pot}
If $a\ne b$, then $m>0$ and $v=(h_b-h_a)/(2m)$ satisfies $Lv=e_a-e_b$.
\end{lemma}
\begin{proof}
$2m=\sum_x\deg x\ge\deg a>0$. Also $Lv=\big((\deg-2m e_b)-(\deg-2m e_a)\big)/(2m)=e_a-e_b$.
\end{proof}

\begin{theorem}[\texttt{commute\_time\_identity}]
For every finite connected simple graph $G$ and all vertices $a,b$:
both tail series converge; a unit-current potential exists; for $a\ne b$,
$(h_b-h_a)/(2m)$ is one; every unit-current potential $v$ gives $\kappa(a,b)=2m\,(v(a)-v(b))$;
$\kappa(a,b)=2m\,\Reff(a,b)$; $\kappa(a,b)=\kappa(b,a)$; and $\Reff(a,b)=\Reff(b,a)$.
\end{theorem}
\begin{proof}
If $a=b$, both sides are $0$. Otherwise, two unit potentials differ by an element of $\ker L$,
i.e.\ by a constant (\texttt{potential\_diff\_eq}). So for any unit potential $v$,
$v(a)-v(b)=\big[(h_b(a)-h_a(a))-(h_b(b)-h_a(b))\big]/(2m)=(H(a,b)+H(b,a))/(2m)$, using
$H(a,a)=H(b,b)=0$. The same holds for $\Reff$ (\texttt{effRes\_eq}). $\kappa$ is symmetric by
definition. $\Reff$ is symmetric because $-v$ is a unit potential from $b$ to $a$
(\texttt{effRes\_comm}).
\end{proof}

\section{Lean formalization and axiom audit}

The file \texttt{lean/Conjecture3419/Basic.lean} (315 lines, \texttt{import Mathlib} only) contains
all the definitions and lemmas above, in namespace \texttt{C3419}. The main theorem is as follows
(Lean's conjunction, existential, inequality, implication and universal quantifier are written here
in ASCII as \verb|/\|, \verb|exists|, \verb|!=|, \verb|->| and \verb|forall|):
\begin{verbatim}
theorem commute_time_identity (hG : G.Connected) (a b : V) :
  (Summable fun t => survProb G b t a) /\
  (Summable fun t => survProb G a t b) /\
  (exists v, IsUnitPotential G a b v) /\
  (a != b -> IsUnitPotential G a b fun x =>
    (hittingTime G x b - hittingTime G x a) / (2 * #G.edgeFinset)) /\
  (forall v, IsUnitPotential G a b v ->
    commuteTime G a b = 2 * #G.edgeFinset * (v a - v b)) /\
  commuteTime G a b = 2 * #G.edgeFinset * effRes G a b /\
  commuteTime G a b = commuteTime G b a /\ effRes G a b = effRes G b a
\end{verbatim}
Here \texttt{V} is a \texttt{Fintype} with decidable equality and \texttt{G : SimpleGraph V}. The
command \texttt{lake build} succeeds. \texttt{\#print axioms C3419.commute\_time\_identity} reports
only \texttt{[propext, Classical.choice, Quot.sound]}. There is no \texttt{sorry} and no
\texttt{native\_decide}.

\end{document}
